\paragraph{Calibration size: the strongest gains occur with limited data.}
At $n=30$, the mean paired volume reductions are 5.64\% (Gaussian),
3.11\% (correlated Gaussian), 7.36\% (Laplace), and 16.44\% (skewed
lognormal). Their 95\% Monte Carlo half-widths are respectively 0.67,
0.34, 1.01, and 2.32 percentage points. At $n=80$ the reductions narrow
to 0.77--1.37\%; at $n=200$ they are 0.23--0.34\%.
The envelope therefore offers a distribution-free direction of improvement,
but its magnitude is strongly setting dependent. The skewed, small-sample
case is the clearest empirical gain in this grid.

Across all 2,160 paired trials and 12,960 coordinate comparisons, there are
no containment violations. In the size study, envelope mean joint coverage
ranges from 0.9006 to 0.9135 and stays close to the 0.9 target. It is slightly
lower than old-shortcut coverage, as nesting requires; shrinking a conservative
region is the intended efficiency gain. The experiment supports the theorem
but does not substitute for it.

\paragraph{Different target levels.}
At $n=80$, paired volume reductions span 1.05--1.64\% for Gaussian noise
and 0.85--1.85\% for Laplace across the four targets. The largest reductions
in both families occur at $\alpha=0.05$, but the curves are not generally
monotone in $\alpha$. Mean envelope coverage exceeds its nominal target by
0.03--1.15 percentage points across these settings. Thus the refinement
persists across targets without a universal claim about how its magnitude
must change with $\alpha$.
